%----------------------------------------------------------------------------------------
% ACADEMIC CURRICULUM VITAE — HAMZA ABBAS
% Mechanical Engineering Lecturer | Thermofluids · Porous Media · Microfluidics
%----------------------------------------------------------------------------------------

\documentclass[11pt,a4paper]{article}
\usepackage[top=1.8cm,bottom=1.8cm,left=2cm,right=2cm]{geometry}
\usepackage{enumitem}
\usepackage{titlesec}
\usepackage{hyperref}
\usepackage{parskip}
\usepackage{tabularx}
\usepackage{xcolor}
\usepackage{fancyhdr}
\usepackage{lastpage}
\usepackage{booktabs}
\usepackage{longtable}
\usepackage[T1]{fontenc}
\usepackage[utf8]{inputenc}
\usepackage{textcomp}
\usepackage{microtype}

\definecolor{accent}{RGB}{0,70,127}

\hypersetup{
  colorlinks=true,
  urlcolor=accent,
  linkcolor=accent
}

\titleformat{\section}{\large\bfseries\color{accent}}{}{0em}{}[\titlerule]
\titlespacing{\section}{0pt}{10pt}{5pt}

\pagestyle{fancy}
\fancyhf{}
\renewcommand{\headrulewidth}{0pt}
\fancyfoot[C]{\small\color{gray}
  Hamza Abbas $\cdot$ CV $\cdot$ Page \thepage\ of \pageref{LastPage}}

\setlist[itemize]{leftmargin=*, noitemsep, topsep=2pt}

%========================================================================================
\begin{document}
%========================================================================================

%--- HEADER ---
\begin{center}
  {\LARGE\bfseries\color{accent} Hamza Abbas}\\[4pt]
  Lecturer, Faculty of Mechanical Engineering $\cdot$ GIK Institute, Topi, KPK,
  Pakistan\\[3pt]
  \href{mailto:hamza.abbas0133@gmail.com}{hamza.abbas0133@gmail.com}
  \quad$|$\quad +92\,348\,0555875
  \quad$|$\quad
  \href{https://linkedin.com/in/hamza-abbas1}{linkedin.com/in/hamza-abbas1}
\end{center}

%--- PROFILE ---
\section{Profile}

Mechanical engineer and lecturer with experience in thermal-fluid sciences, heat
transfer, fluid mechanics, and numerical modelling. My academic and project work
combines transport analysis, simulation, and experimental validation across
porous-media, microfluidic, and energy-system applications. I use COMSOL
Multiphysics for flow and heat-transfer modelling, and Aspen Plus for
thermodynamic cycle simulation. My MS thesis focused on fluid transport in
heterogeneous porous assemblies, where I derived and experimentally validated a
modified Lucas--Washburn equation with below 8\% error. I am particularly
interested in doctoral research involving heat transfer, turbulent flows, numerical
optimisation, and thermofluid system design.

%--- WORK EXPERIENCE ---
\section{Work Experience}

\textbf{University Lecturer} -- GIK Institute, Faculty of Mechanical Engineering
\hfill \textit{Sep 2024 -- Present}\\
\textit{Topi, Pakistan}
\begin{itemize}
  \item Teach HVAC (ME439), Engineering Dynamics (ME212), and CAD/CAM Lab
        (ME418-L) to undergraduate students; coordinate Thermofluids-I Lab
        (ME346-L) and manage the Fluid Mechanics Laboratory.
  \item Serve as \textbf{Social Media Coordinator, FME} -- manage the faculty's
        official Facebook and LinkedIn pages, promoting research output, events,
        and student achievements.
  \item Act as \textbf{Batch Advisor} for Batch~34 (Non-Lead role), providing
        academic guidance and pastoral support.
  \item Support Outcome-Based Education mapping for courses aligned with
        Washington Accord standards; contributed to writing Criterion~0 of the
        Self-Assessment Report (SAR) for institutional accreditation.
  \item Introduced flipped-classroom and active-learning methods in ME439; revised
        course title from ``Refrigeration \& Air-Conditioning'' to ``Heating,
        Ventilation and Air-Conditioning'' with updated content including
        refrigerant selection.
  \item Initiated replacement of FME and Management elective courses with
        accredited Coursera/edX equivalents to broaden student learning options.
  \item \textbf{Lead Organiser} -- OGDCL Technical Workshop (Oil \& Gas
        Development Company Limited); coordinated all logistical, technical, and
        communication aspects of the industry-academia workshop hosted at GIK
        Institute (Summer~2025).
  \item \textbf{Co-Organiser} -- 4th International Conference on Modern
        Technologies in Mechanical, Material, and Chemical Engineering
        (\textbf{MTME~2025}), GIK Institute, Topi, Pakistan (Apr~2025).
  \item \textbf{Co-Organiser} -- 4th International Conference on Modern
        Technologies in Mechanical, Material, and Chemical Engineering
        (\textbf{MTME~2026}), GIK Institute, Topi, Pakistan.
  \item \textbf{Organiser} -- 6th International Workshop on Functional Reverse
        Engineering of Machine Tools (\textbf{WRE~2024}), GIK Institute,
        Topi, Pakistan.
  \item \textbf{Led Industrial Trip} to KSB Pumps (Dec~2025) for undergraduate
        students; arranged industrial visit to DG Khan Cement Company (Fall~2024);
        facilitated student visit to OGDCL Head Office, Islamabad (Fall~2024).
  \item Served on Faculty Study Board (2024, 2025), CTD Committee (2024), and
        Industrial Advisory Board (2024); recorded minutes and proposed course
        title revisions.
\end{itemize}

\vspace{4pt}
\textbf{Graduate Research Assistant} -- Biofuels \& Microfluidics Lab, GIK
Institute \hfill \textit{Feb 2022 -- May 2024}\\
\textit{Topi, Pakistan}
\begin{itemize}
  \item Used COMSOL Multiphysics to model flow behaviour, porous-media transport,
        electric fields, and heat transfer in microfluidic and porous structures.
  \item Derived a modified Lucas--Washburn equation for wetting-front propagation
        in heterogeneous porous assemblies and validated it experimentally to
        below 8\% error.
  \item Investigated capillary flow dynamics in paper-based multilayer systems as
        a function of substrate thickness, contact area, and fluid volume.
  \item Fabricated paper-based and digital microfluidic devices using CO$_2$ laser
        micromachining, photolithographic electrode patterning on PET, and thermal
        lamination.
  \item Used ImageJ for experimental image analysis and OriginPro for curve
        fitting and data analysis.
  \item Supported Thermo-Fluid Lab~1, Heat Transfer, Measurement \&
        Instrumentation, and Micro \& Nano Fabrication courses.
\end{itemize}

\vspace{4pt}
\textbf{Project Research Assistant} -- HEC Project TTSF-74, GIK Institute
\hfill \textit{Feb 2022 -- Feb 2024}\\
\textit{Topi, Pakistan}
\begin{itemize}
  \item Contributed to an HEC-funded project developing a flexible
        time-temperature indicator for food quality control; supported early
        business-model development for the spin-off startup ``Flow Time.''
  \item Co-inventor: patent application submitted for \textit{``An enclosed
        fan-shaped flexible time-temperature indicator''} (local filing).
\end{itemize}

\vspace{4pt}
\textbf{Mechanical Engineering Intern} -- DG Khan Cement Company Limited
\hfill \textit{Aug 2020 -- Sep 2020}\\
\textit{Chakwal, Pakistan}
\begin{itemize}
  \item Studied plant schematics, joined maintenance crews, and gained hands-on
        exposure to heavy industrial machinery in a cement manufacturing facility.
\end{itemize}

\vspace{4pt}
\textbf{Mechanical Engineering Intern} -- KKB Engineering
\hfill \textit{Mar 2020 -- Apr 2020}\\
\textit{Chakwal, Pakistan}
\begin{itemize}
  \item Observed maintenance of a limestone crusher, raw mill, and kiln; spent
        time in the quality control department.
\end{itemize}

%--- EDUCATION ---
\section{Education}

\textbf{M.S. in Mechanical Engineering} -- Thermofluid Specialisation, GIK
Institute \hfill \textit{Feb 2022 -- Jun 2024}\\
\textit{Topi, Pakistan}
\begin{itemize}
  \item GPA: 3.67/4.00. Awarded Merit-Based Graduate Scholarship and Research
        Assistantship (GA-1).
  \item Thesis: \textit{Dynamics of Fluid Flow Between Two Porous Membranes.}
        Derived and experimentally validated a modified Lucas--Washburn equation
        for stratified porous assemblies with dissimilar permeabilities.
  \item Relevant coursework: Intermediate Fluid Mechanics (A), Mass/Momentum/Energy
        Transport (A-), Energy Management \& Auditing (A), Microfluidic Systems
        (A-), Advanced Stress Analysis (B), Elements of Vibration \& Feedback
        Control (B+).
\end{itemize}

\vspace{4pt}
\textbf{B.S. in Mechanical Engineering} -- GIK Institute
\hfill \textit{Aug 2017 -- Jul 2021}\\
\textit{Topi, Pakistan}
\begin{itemize}
  \item Senior Design Project: designed and simulated an integrated ORC+VARS
        waste heat recovery system for DG Khan Cement using Aspen Plus with
        Peng--Robinson EOS; achieved 17.9\% combined efficiency and 8.7\,MW net
        output using n-Hexane as the working fluid.
\end{itemize}

%--- RESEARCH & INDUSTRIAL PROJECTS ---
\section{Research \& Industrial Projects}

\textbf{Industrial Water Heating System -- Pakistan Tobacco Company (Co-I)}
\hfill \textit{Apr 2025 -- Present}\\
Co-investigator on an industry-funded project designing a large-scale water
heating system (600\,kW) integrated with grid-tied photovoltaic generation;
leading technical specification and cost reporting.

\vspace{4pt}
\textbf{Cooling and Heating Load Calculation -- Independent Research}
\hfill \textit{Dec 2025 -- Feb 2026}\\
Conducted a detailed building thermal-load analysis (cooling and heating) as
independent research; output: one conference paper (in preparation).

\vspace{4pt}
\textbf{ORC Feasibility Study -- Low-Temperature Industrial Heat Source}
\hfill \textit{Feb -- Mar 2025}\\
Assessed the techno-economic viability of an Organic Rankine Cycle on a low-grade
industrial waste heat stream, including working fluid selection, cycle layout, and
economic payback.

\vspace{4pt}
\textbf{ORC+VARS Waste Heat Recovery -- DG Khan Cement}
\hfill \textit{Sep 2020 -- Jul 2021}\\
Modelled an integrated ORC+VARS system in Aspen Plus for kiln waste heat
recovery. Final design: 8.7\,MW net output and 17.9\% combined efficiency.

\vspace{4pt}
\textbf{Packed-Bed Latent Heat Storage -- COMSOL Simulation}
\hfill \textit{Feb -- Jun 2020}\\
Simulated transient thermal behaviour in a packed-bed PCM storage tank using
COMSOL Multiphysics.

%--- PUBLICATIONS ---
\section{Publications}

\textit{Journal Articles -- Under Review}
\begin{itemize}
  \item Ali, M.; Jafry, A.T.; \textbf{Abbas, H.}; Lee, M.; Abbasi, M.S.; Shah,
        S.F.; Ajab, H.; Kamran, M.; Jafry, A.H. ``Revealing the Mechanics of
        Highly Resistive Fluid Flow in a Leaf Architecture using Paper
        Microfluidics.''
        \textbf{\textit{Journal of Fluid Mechanics}} (Cambridge), under review,
        2025.
  \item \textbf{Abbas, H.}; Jafry, A.T.; Ali, M.; Ajab, H.; Cheema, T.A.
        ``Impact of paper thickness, contact area, and fluid volume on flow
        dynamics in paper-based microfluidics.''
        \textbf{\textit{Langmuir}} (ACS), under review, 2025.
\end{itemize}

\vspace{4pt}
\textit{Conference Proceedings -- Published}
\begin{itemize}
  \item Ullah, H.; Ali, M.; Naeem, N.; Iqbal, F.; \textbf{Abbas, H.}; Jafry,
        A.T.; Ajab, H. ``Droplet Actuation Enhancement through Voltage Control
        and Hydrophobic Coating Selection.''
        \textbf{\textit{MATEC Web of Conferences}}, vol.~398, p.~01020, EDP
        Sciences, 2024.
  \item \textbf{Abbas, H.}; Ali, M.; Ullah, H.; Jafry, A.T. ``Actuation and
        Flow Control in Paper-Based Microfluidics by Varying Thickness of Storage
        Reservoir.''
        \textbf{\textit{MATEC Web of Conferences}}, vol.~381, p.~01010, EDP
        Sciences, 2023.
  \item Ali, M.; Jafry, A.T.; \textbf{Abbas, H.}; Ullah, H. ``Flow Control in
        Paper-Based Microfluidic Channel by Using Hydrophobic Tape Barrier on the
        Paper Surface.''
        \textbf{\textit{MATEC Web of Conferences}}, vol.~381, p.~01009, EDP
        Sciences, 2023.
  \item \textbf{Abbas, H.}; Ali, M.; Naeem, N.; Ullah, H.; Jafry, A.T. ``Effect
        of Storage Reservoir on Fluid Velocity in Lateral Flow Paper Device.''
        \textbf{\textit{Engineering Proceedings}} (MDPI), vol.~45(1), p.~5, 2023.
  \item Ali, M.; Ullah, H.; Ali, M.; Naeem, N.; \textbf{Abbas, H.}; Jafry, A.T.
        ``Dynamic Contact Angle Variation with Applied Voltage and Droplet Volume
        in Digital Microfluidics.''
        \textbf{\textit{Engineering Proceedings}} (MDPI), vol.~45(1), p.~31,
        2023.
  \item Ullah, H.; Shoaib, F.; Zahid, S.D.; Mahmood, M.I.; Ali, M.; Ali, M.;
        \textbf{Abbas, H.}; Jafry, A.T. ``Design and Fabrication of Digital
        Microfluidics Device for Lab-on-a-Chip Applications.''
        \textbf{\textit{Engineering Proceedings}} (MDPI), vol.~45(1), p.~7, 2023.
  \item Ali, M.; Hussain, G.; \textbf{Abbas, H.}; Ullah, H.; Jafry, A.T. ``Flow
        Control in Paper-Based Microfluidics using Variable Porosity Channel.''
        \textbf{\textit{Engineering Proceedings}} (MDPI), vol.~23(1), p.~33,
        2022.
\end{itemize}

%--- PATENT ---
\section{Patent}

\begin{itemize}
  \item \textbf{Abbas, H.}; et al. ``An Enclosed Fan-Shaped Flexible
        Time-Temperature Indicator.'' Patent application submitted (local filing,
        Pakistan), 2025. \textit{(In submission)}
\end{itemize}

%--- CONFERENCE & WORKSHOP ORGANISATION ---
\section{Conference \& Workshop Organisation}

\begin{itemize}
  \item \textbf{Co-Organiser} -- 4th International Conference on Modern
        Technologies in Mechanical, Material, and Chemical Engineering
        (\textbf{MTME~2026}), GIK Institute, Topi, Pakistan.
  \item \textbf{Co-Organiser} -- 4th International Conference on Modern
        Technologies in Mechanical, Material, and Chemical Engineering
        (\textbf{MTME~2025}), GIK Institute, Topi, Pakistan (Apr~2025).
  \item \textbf{Lead Organiser} -- OGDCL Technical Workshop (Oil \& Gas
        Development Company Limited), industry-academia workshop, GIK Institute,
        Topi, Pakistan (Summer~2025).
  \item \textbf{Organiser} -- 6th International Workshop on Functional Reverse
        Engineering of Machine Tools (\textbf{WRE~2024}), GIK Institute,
        Topi, Pakistan.
\end{itemize}

%--- TECHNICAL SKILLS ---
\section{Technical Skills}

\begin{tabularx}{\textwidth}{@{}>{\bfseries}lX@{}}
Simulation \& Modelling
  & COMSOL Multiphysics (flow, heat transfer, porous media, electric fields);
    Aspen Plus (thermodynamic cycle modelling, process simulation) \\[3pt]
Numerical / Analysis Tools
  & MATLAB, Python, OriginPro, ImageJ \\[3pt]
Experimental Fabrication
  & CO$_2$ laser micromachining, thermal lamination, photolithographic electrode
    patterning on PET \\[3pt]
Other
  & Microsoft Office; Moodle LMS; OBE/SAR documentation; English C1 (CEFR) \\
\end{tabularx}

%--- MEMBERSHIPS & AWARDS ---
\section{Professional Memberships \& Awards}
\begin{itemize}
  \item Registered Mechanical Engineer, Pakistan Engineering Council
        (Lifetime member) -- Sep 2021.
  \item Merit-Based Graduate Scholarship \& Assistantship (GA-1), GIK Institute
        -- Feb 2022.
  \item 5th place, DICE AFS-2022 Innovation Competition, MNS University of
        Agriculture, Multan -- Oct 2022.
\end{itemize}

%--- REFERENCES ---
\section{References}

\begin{tabularx}{\textwidth}{@{}XX@{}}
\textbf{Dr Ali Turab Jafry}       & \textbf{Dr Taqi Ahmad Cheema} \\
Associate Professor, FME          & Dean / Associate Professor, FME \\
GIK Institute, Topi               & GIK Institute, Topi \\
\href{mailto:ali.turab@giki.edu.pk}{ali.turab@giki.edu.pk}
  & \href{mailto:tacheema@giki.edu.pk}{tacheema@giki.edu.pk} \\
\end{tabularx}

%========================================================================================
% PAGE 2 — PhD APPLICATION TRACKER
%========================================================================================
\newpage

\begin{center}
  {\LARGE\bfseries\color{accent} PhD Application Tracker}\\[4pt]
  {\large Hamza Abbas $\cdot$ 2025/2026 Cycle}\\[2pt]
  \small Updated: May 2026
\end{center}

\vspace{8pt}
\small
\renewcommand{\arraystretch}{1.65}

\begin{longtable}{@{} p{3.6cm} p{3.0cm} p{2.8cm} p{1.9cm} p{3.2cm} @{}}
\toprule
\textbf{Institution \& Lab} & \textbf{PI} & \textbf{Focus / Pillar}
  & \textbf{Deadline} & \textbf{Status / Action} \\
\midrule
\endhead

%--- Target 1 ---
\textbf{Queensland University of Technology (QUT)}
\newline Brisbane, Australia
\newline \textit{Multiphase Flow \& Porous Media Group}
&
Dr.\ Zhongzheng Wang \newline
Prof.\ Emilie Sauret
&
Structural capillarity; engineered porous media; LBM; wicking
\newline \textit{Pillar 2}
&
\textcolor{accent}{\textbf{31 Jul 2026}}
&
\textcolor{accent}{\textbf{OPEN}} \newline
Fully funded. Email supervisor first; EOI after interview. \\

\midrule

%--- Target 2 ---
\textbf{University of Stuttgart}
\newline Stuttgart, Germany
\newline \textit{SFB 1313 -- Porous Media}
&
Dr.\ Nikolaos Karadimitriou
&
Multiphase flow; capillary pressure; interfacial phenomena
\newline \textit{Pillar 2}
&
\textcolor{red}{\textbf{15 Apr 2026}}
&
\textcolor{red}{\textbf{PASSED}} \newline
Monitor for next cycle. \\

\midrule

%--- Target 3 ---
\textbf{University of Strathclyde}
\newline Glasgow, UK
\newline \textit{Mech.\ \& Aerospace Eng.}
&
[Confirm PI on FindAPhD listing]
&
Microfluidic fabrication; diagnostic flow analysis
\newline \textit{Pillar 3}
&
\textcolor{accent}{\textbf{31 Jul 2026}}
&
\textcolor{accent}{\textbf{OPEN}} \newline
Formal application via FindAPhD. \\

\midrule

%--- Target 4 ---
\textbf{University of Toronto}
\newline Toronto, Canada
\newline \textit{Wheeler Microfluidics Lab}
&
Prof.\ Aaron Wheeler \newline
\href{mailto:aaron.wheeler@utoronto.ca}{\footnotesize aaron.wheeler@utoronto.ca}
&
Digital microfluidics; EWOD; paper diagnostics
\newline \textit{Pillar 3}
&
\textcolor{orange}{\textbf{31 May 2026}} \newline
{\tiny BME Round 2}
&
\textcolor{orange}{\textbf{URGENT}} \newline
Cold email + BME application now. \\

\midrule

%--- Target 5 ---
\textbf{UC Merced}
\newline California, USA
\newline \textit{Aghilinejad Group}
&
Prof.\ Arian Aghilinejad \newline
\href{mailto:aaghilinejad@ucmerced.edu}{\footnotesize aaghilinejad@ucmerced.edu}
&
Biofluid mechanics; microfluidics; ML diagnostics
\newline \textit{Pillar 3}
&
\textcolor{red}{\textbf{15 Jan 2026}} \newline
{\tiny Fall 2026 passed} \newline
\textit{$\sim$Dec 2026} \newline
{\tiny Fall 2027 entry}
&
\textcolor{red}{Jan 2026 passed.} \newline
Cold email now; \newline apply Dec 2026. \\

\midrule

%--- Target 6 ---
\textbf{University of Washington}
\newline Seattle, USA
\newline \textit{Lutz Lab -- Bioengineering}
&
Prof.\ Barry Lutz \newline
\href{mailto:blutz@uw.edu}{blutz@uw.edu}
&
Capillary paper microfluidics; point-of-care diagnostics
\newline \textit{Pillars 2 + 3}
&
\textit{$\sim$Dec 2026} \newline
{\tiny Fall 2027 entry}
&
\textcolor{red}{Dec 2025 passed.} \newline
Email now; apply Dec 2026. \\

\midrule

%--- Target 7 ---
\textbf{UCLA}
\newline Los Angeles, USA
\newline \textit{Micro \& Nano Mfg.\ Lab}
&
Prof.\ Chang-Jin ``CJ'' Kim
&
EWOD; digital microfluidics; droplet actuation
\newline \textit{Pillar 3}
&
\textit{$\sim$Dec 2026} \newline
{\tiny Fall 2027 entry}
&
\textcolor{red}{Dec 2025 passed.} \newline
Email now; apply Dec 2026. \\

\midrule

%--- Target 8 ---
\textbf{University of Birmingham}
\newline Birmingham, UK
\newline \textit{Centre for Energy Systems}
&
[Identify via EPSRC award search]
&
ORC; waste heat recovery; Aspen Plus; thermal systems
\newline \textit{Pillar 1}
&
Rolling / EPSRC CDT
&
Monitor FindAPhD. \newline
Confirm intl.\ fee eligibility. \\

\bottomrule
\end{longtable}

\vspace{6pt}
\noindent
\textbf{Pillar key:} \
1 = ORC/Waste-Heat Simulation \quad
2 = Porous Media \& Capillary Flow \quad
3 = Microfluidics Fabrication \& Diagnostics

\vspace{8pt}
\noindent
\textbf{Deadline key:} \
\textcolor{accent}{$\blacksquare$}~Open \quad
\textcolor{orange}{$\blacksquare$}~Urgent \quad
\textcolor{red}{$\blacksquare$}~Closed/Next cycle

\vspace{10pt}
\noindent\textit{Always verify deadlines directly on the institution's official
scholarship or admissions page before applying. Last updated: May 2026.}

\end{document}
